\chapter*{Resumo}

Os jogos digitais multiplayer oferecem ambientes de interação que podem contribuir para o desenvolvimento de habilidades sociais. Este trabalho teve como objetivo investigar o potencial dos jogos digitais como ferramentas de treinamento dessas habilidades. Para isso, foram realizadas revisões da literatura e um experimento utilizando o jogo Among’Us, com 15 participantes divididos entre um grupo experimental (n = 9) e um grupo controle (n = 6). As habilidades sociais foram avaliadas antes e após o experimento por meio do Inventário de Habilidades Sociais (IHS-Del-Prette), complementado por questionários sobre as percepções dos participantes. Os dados foram analisados utilizando os testes de Wilcoxon e Mann–Whitney, com correção de Holm. Os resultados apresentaram aumentos médios em três dos cinco fatores avaliados no grupo experimental, porém sem diferenças estatisticamente significativas. Apesar disso, todos os participantes do grupo experimental relataram perceber melhorias em suas habilidades sociais e conhecer novas pessoas. Conclui-se que, embora não tenha sido possível comprovar a eficácia da intervenção, os resultados exploratórios indicam o potencial dos jogos multiplayer comerciais como ambientes de interação social, destacando a necessidade de estudos com amostras maiores.
